\documentclass[../../main.tex]{subfiles}

\begin{document}

\section{半线性波动方程}\label{section:ch02-4-波动方程}

现在研究具非线性项 $f(u)$ 的波动方程。我们从一个情形开始，其中解在时间上全局存在。

\subsection{符号条件}\label{subsection:ch12-符号条件}

考虑问题
\begin{align}
\begin{cases}
u_{tt}-\Delta u+f(u)=0 & \text{在 } \mathbb{R}^n\times(0,\infty) \text{ 中},\\
u=g,\;u_t=h & \text{在 } \mathbb{R}^n\times\{t=0\} \text{ 上}.
\end{cases}
\label{align:ch12-半线性初值}
\end{align}

\begin{theorem}[符号条件下的全局存在性。]\label{theorem:ch12-符号条件全局存在}设 $f$ 光滑，$f(0)=0$，且
\begin{align}
zf(z)\geqslant 0 \quad (z\in\mathbb{R}).
\label{align:ch12-符号条件}
\end{align}
假设 $g,h\in C_c^\infty(\mathbb{R}^n)$。则存在初值问题\eqref{align:ch12-半线性初值} 的唯一全局光滑解。
\end{theorem}

\begin{proof}
1. 由 \S12.2\cref{theorem:ch12-短时间存在性} 存在短时间解。设 $T^*>0$ 为此解存在的最大时刻。我们需要证明若 $T^*<\infty$，则可将解延拓到时间 $T^*$ 之后，矛盾。

2. 首先假设 $T^*<\infty$，并断言
\begin{align}
\|u\|_{H^k(\mathbb{R}^n)}+\|u_t\|_{H^{k-1}(\mathbb{R}^n)}\leqslant C \quad (0\leqslant t<T^*)
\label{align:ch12-Hk界}
\end{align}
对某常数 $C=C(T^*)$ 和某 $k>n/2+1$。一旦此成立，\S12.2\cref{theorem:ch12-短时间存在性} 的证明表明解可延拓超过 $T^*$。

3. 我们首先验证
\begin{align}
\|u\|_{L^\infty(\mathbb{R}^n\times[0,T^*))}+\|u_t\|_{L^\infty(\mathbb{R}^n\times[0,T^*))}\leqslant C.
\label{align:ch12-Linfty界}
\end{align}
为此注意\eqref{align:ch12-符号条件} 蕴含 $F(u)=\int_0^uf(w)\,\mathrm{d}w\geqslant 0$。因此能量守恒（\S12.1\cref{theorem:ch12-能量守恒}）给出
\begin{align}
\int_{\mathbb{R}^n}\frac{1}{2}(u_t^2+|Du|^2)+F(u)\,\mathrm{d}x=E(0)
\label{align:ch12-能量守恒2}
\end{align}
对所有 $t\geqslant 0$。

接下来建立 $L^\infty$ 估计。回忆线性波动方程的表示公式。设 $v$ 是
\begin{align}
\begin{cases}
v_{tt}-\Delta v=0 & \text{在 } \mathbb{R}^n\times(0,\infty) \text{ 中},\\
v=g,\;v_t=h & \text{在 } \mathbb{R}^n\times\{t=0\} \text{ 上}.
\end{cases}
\label{align:ch12-线性波动初值}
\end{align}
Duhamel 原理给出
\begin{align}
u(x,t)=v(x,t)-\int_0^t\frac{\sin((t-s)\sqrt{-\Delta})}{\sqrt{-\Delta}}f(u)(\cdot,s)(x)\,\mathrm{d}s.
\label{align:ch12-Duhamel}
\end{align}
现在 $v$ 是光滑有界的。此外由有限传播速度和能量守恒，$u$ 对每个时刻具有空间紧支集。

4. 对每个 $\alpha$ 多重指标，$D^\alpha u$ 满足
\begin{align*}
(D^\alpha u)_{tt}-\Delta(D^\alpha u)+f'(u)D^\alpha u=R_\alpha(Du,\cdots,D^{|\alpha|}u),
\end{align*}
其中 $R_\alpha$ 是光滑函数。由于 $u$ 有界，$f'(u)$ 有界。能量估计和 Gronwall 不等式给出 $D^\alpha u$ 的 $L^2$ 界。归纳推出所有导数在 $H^k$ 中有界，即\eqref{align:ch12-Hk界}。
\end{proof}

\subsection{三维空间}\label{subsection:ch12-三维空间}

现在考虑非线性项不取非负性条件\eqref{align:ch12-符号条件} 的情形。我们专注于 $n=3$ 空间维数。

\begin{theorem}[三维中的短时间存在性与 $L^\infty$ 爆破。]\label{theorem:ch12-三维短时间爆破}考虑
\begin{align}
\begin{cases}
u_{tt}-\Delta u+f(u)=0 & \text{在 } \mathbb{R}^3\times(0,\infty) \text{ 中},\\
u=g,\;u_t=h & \text{在 } \mathbb{R}^3\times\{t=0\} \text{ 上},
\end{cases}
\label{align:ch12-三维初值}
\end{align}
其中 $f\in C^\infty(\mathbb{R})$，$f(0)=0$，$g,h\in C_c^\infty(\mathbb{R}^3)$。
\begin{enumerate}[(i)]
\item 存在唯一光滑解 $u$ 存在于 $\mathbb{R}^3\times[0,T)$，其中 $T>0$ 是最大的。
\item 若 $T<\infty$，则
\begin{align}
\limsup_{t\to T}\|u(\cdot,t)\|_{L^\infty(\mathbb{R}^3)}=\infty.
\label{align:ch12-爆破}
\end{align}
\end{enumerate}
\end{theorem}

\begin{proof}
(i) 由 \S12.2\cref{theorem:ch12-短时间存在性} 得。

(ii) 假设 $T<\infty$ 但\eqref{align:ch12-爆破} 不成立，则 $\|u\|_{L^\infty(\mathbb{R}^3\times[0,T))}\leqslant C$。回忆 Kirchhoff 公式（\S2.4.2），可写
\begin{align}
u(x,t)=v(x,t)-\frac{1}{4\pi}\int_{B(x,t)}\frac{f(u^*)}{|x-y|}\,\mathrm{d}y,
\label{align:ch12-Kirchhoff}
\end{align}
其中 $v$ 是齐次线性波动方程的解，$u^*(y)\coloneqq u(y,t-|y-x|)$。由于 $g,h$ 光滑，$v$ 有界。因此
\begin{align*}
|u(x,t)|\leqslant C+C\int_{B(x,t)}\frac{|f(u^*)|}{|x-y|}\,\mathrm{d}y.
\end{align*}
因 $u$ 有界且 $f$ 光滑，$|f(u^*)|\leqslant C|u^*|$。从而
\begin{align*}
\|u(\cdot,t)\|_{L^\infty}\leqslant C+C\int_0^t\|u(\cdot,s)\|_{L^\infty}\,\mathrm{d}s.
\end{align*}
Gronwall 不等式给出 $u$ 在 $\mathbb{R}^3\times[0,T]$ 上有界。然后如\cref{theorem:ch12-符号条件全局存在} 证明中估计所有导数，可将解延拓超过 $T$，矛盾。
\end{proof}

\subsection{次临界幂非线性}\label{subsection:ch12-次临界幂非线性}

现在转向三维中具幂型非线性的非线性波动方程：
\begin{align}
\begin{cases}
u_{tt}-\Delta u+|u|^{p-1}u=0 & \text{在 } \mathbb{R}^3\times(0,\infty) \text{ 中},\\
u=g,\;u_t=h & \text{在 } \mathbb{R}^3\times\{t=0\} \text{ 上},
\end{cases}
\label{align:ch12-幂非线性}
\end{align}
能量为
\begin{align*}
E(t)\coloneqq\int_{\mathbb{R}^3}\frac{1}{2}(u_t^2+|Du|^2)+\frac{|u|^{p+1}}{p+1}\,\mathrm{d}y.
\end{align*}
假设 $g,h\in C_c^\infty(\mathbb{R}^3)$。目标是对 $1\leqslant p<5$ 构造全局解。

\begin{theorem}[次临界存在定理。]\label{theorem:ch12-次临界存在}若 $1<p<5$，则存在\eqref{align:ch12-幂非线性} 的唯一全局光滑解。
\end{theorem}

\S12.4 将引入更精细的技术处理临界情形 $p=5$。

\begin{proof}
1. 由\cref{theorem:ch12-三维短时间爆破}，可假设已有光滑解 $u$ 存在于 $\mathbb{R}^3\times(0,T)$ 某 $T>0$。需证明
\begin{align*}
\|u\|_{L^\infty(\mathbb{R}^3\times(0,T))}\leqslant C
\end{align*}
对某常数 $C=C(T)$。

设 $v$ 是齐次线性波动方程的解，则如前一证明
\begin{align*}
u(x,t)=v(x,t)-\frac{1}{4\pi}\int_{B(x,t)}\frac{|u^*|^{p-1}u^*}{|x-y|}\,\mathrm{d}y.
\end{align*}
因 $g,h$ 光滑，$v$ 在 $\mathbb{R}^3\times[0,T]$ 上有界。因此
\begin{align}
|u(x,t)|\leqslant C+I(x,t) \quad (x\in\mathbb{R}^3,\;0\leqslant t\leqslant T)
\label{align:ch12-u界}
\end{align}
其中
\begin{align*}
I(x,t)\coloneqq\int_{B(x,t)}\frac{|u^*|^p}{|x-y|}\,\mathrm{d}y.
\end{align*}

2. 然后
\begin{align}
I(x,t)\leqslant\left(\int_{B(x,t)}\frac{|u^*|^2}{|x-y|^2}\,\mathrm{d}y\right)^{1/2}\left(\int_{B(x,t)}|u^*|^{2(p-1)}\,\mathrm{d}y\right)^{1/2}.
\label{align:ch12-I估计}
\end{align}
Hardy 不等式（\S5.8.4\cref{theorem:ch05-Hardy}）蕴含
\begin{align*}
\int_{B(x,t)}\frac{|u^*|^2}{|y-x|^2}\,\mathrm{d}y\leqslant C\int_{B(x,t)}\frac{|Du^*|^2+(u^*)^2}{t^2}\,\mathrm{d}y.
\end{align*}
现在 $Du^*=Du-\nu u_t$，$\nu=\dfrac{y-x}{|y-x|}$。能量通量恒等式\eqref{align:ch12-通量恒等式} 蕴含
\begin{align*}
\int_{B(x,t)}|Du^*|^2\,\mathrm{d}y=\frac{1}{\sqrt{2}}\int_{\Gamma(x,t)}|u_t\nu-Du|^2\,\mathrm{d}S\leqslant E(0).
\end{align*}
回忆\cref{theorem:ch05-Poincare}，
\begin{align*}
\int_{B(x,t)}|u^*-(u^*)_{x,t}|^2\,\mathrm{d}y\leqslant Ct^2\int_{B(x,t)}|Du^*|^2\,\mathrm{d}y,
\end{align*}
平均 $(u^*)_{x,t}\coloneqq\fint_{B(x,t)}u^*\,\mathrm{d}y$。因此
\begin{align}
\frac{1}{t^2}\int_{B(x,t)}(u^*)^2\,\mathrm{d}y\leqslant C\int_{B(x,t)}|Du^*|^2\,\mathrm{d}y+Ct|(u^*)_{x,t}|^2.
\label{align:ch12-Poincare}
\end{align}

3. 计算
\begin{align*}
|(u^*)_{x,t}|\leqslant\frac{C}{t^3}\int_{\Gamma(x,t)}|u|\,\mathrm{d}S\leqslant\frac{C}{t^3}\left(\int_{K(x,t)}|u_t|^2\,\mathrm{d}y\mathrm{d}s\right)^{1/2}|K(x,t)|^{1/2}+C\|g\|_{L^\infty}\leqslant Ct^{-1/2}E(0)^{1/2}+C\|g\|_{L^\infty}.
\end{align*}
此估计结合\eqref{align:ch12-I估计}--\eqref{align:ch12-Poincare} 给出
\begin{align}
I(x,t)\leqslant C\left(\int_{B(x,t)}|u^*|^{2(p-1)}\,\mathrm{d}y\right)^{1/2}.
\label{align:ch12-I最终}
\end{align}

4. 假设 $1<p\leqslant 4$。则 $2(p-1)\leqslant 6=2^*$，Sobolev 不等式（\S5.6.3）蕴含
\begin{align*}
\|(u^*)^{p-1}\|_{L^2(B(x,t))}\leqslant C\|u^*\|_{H^1(B(x,t))}^{p-1}\leqslant C.
\end{align*}
因此\eqref{align:ch12-u界} 和\eqref{align:ch12-I最终} 给出 $|u(x,t)|\leqslant C$。

下一种情形 $4<p<5$。则 $2(p-1)=2(p-4)+6=2(p-4)+2^*$，前述计算表明
\begin{align*}
I(x,t)\leqslant C\left(\int_{B(0,t)}|u^*|^{2(p-1)}\,\mathrm{d}y\right)^{1/2}\leqslant C\|u\|_{L^\infty}^{p-4}.
\end{align*}
因此\eqref{align:ch12-u界} 蕴含
\begin{align*}
|u(x,t)|\leqslant C+C\|u\|_{L^\infty}^{p-4}\leqslant C+\frac{1}{2}\|u\|_{L^\infty},
\end{align*}
最后一不等式成立因 $0<p-4<1$。这再次给出 $\|u\|_{L^\infty(\mathbb{R}^3\times(0,T))}$ 的界。
\end{proof}

我们将在\cref{exercise:ch12-小能量p5}中看到，证明可修改以表明对临界幂 $p=5$，只要能量 $E(0)$ 足够小，就存在全局光滑解。下一节引入更先进技术，即使 $p=5$ 时能量大也构造解。

\end{document}